\section{Introducción: la segunda mitad de la IA pasa de la capacidad del modelo a la capacidad del sistema}

Esta sección establece primero las coordenadas de todo el episodio. Yao Shunyu (姚顺雨) publicó «The Second Half» en abril de 2025, donde sostiene que el hilo principal de la IA ya entró en su segunda mitad. La línea principal de la primera mitad se parece más a la capacidad del modelo: preentrenamiento, escalamiento, conversación y capacidades generales; la segunda mitad se parece más a la capacidad del sistema: Agent, entornos de tareas, mecanismos de recompensa, llamadas a herramientas, interfaces de interacción y estructura organizacional. La entrevista parte de la experiencia personal, pero su núcleo real es una pregunta: cuando el modelo es lo bastante potente, ¿cómo entra la inteligencia en el mundo y da lugar a acción, organización y nuevas fronteras?

\begin{figure}[H]
\centering
\includegraphics[width=0.92\textwidth]{second-half-map.png}
\caption{Mapa de la segunda mitad de la IA: de la capacidad del modelo hacia Agent, sistema, interacción y el panorama humano global. Diagrama conceptual propio, elaborado a partir de la conversación de 00:01:45--02:31:20.}
\end{figure}

\begin{knowledgebox}{Lectura de la figura: la segunda mitad no significa «el modelo ya no importa»}
El modelo sigue siendo la base, pero el centro de la discusión pasa de la puntuación de un modelo individual a cómo actúa el sistema: si la tarea es real, si el entorno es lo bastante amplio, si la reward puede definirse, si la interface abre nuevas oportunidades y si el mundo final será unipolar o plural.
\end{knowledgebox}

\begin{importantbox}{Tesis central del episodio}
Un Agent no es un botón de función de un modelo conversacional, sino un paradigma de sistema: debe observar, razonar, actuar y recibir retroalimentación dentro de un entorno, y ampliar continuamente la frontera de las tareas que puede completar mediante herramientas, recompensas y estructuras multiagente.
\end{importantbox}

\begin{knowledgebox}{Tabla general de conceptos clave}
\begin{tabularx}{\textwidth}{p{0.20\textwidth}p{0.34\textwidth}X}
\toprule
Concepto & Significado en este episodio & Malentendido que conviene evitar al estudiar \\
\midrule
Agent & Sistema que interactúa con un entorno e intenta optimizar un objetivo & No es tan simple como «una ventana de chat con algunas herramientas». \\
Environment & El mundo donde el Agent actúa y recibe retroalimentación & Si el entorno es demasiado pequeño, la capacidad no puede extrapolarse. \\
Reward & Señal que evalúa las consecuencias de una acción & La reward no equivale a todo el valor que las personas realmente quieren. \\
Affordance & Posibilidades de acción que el entorno ofrece a quien actúa & El code/API es la mano de la IA, no solo un formato de salida. \\
Interface & Forma en que el usuario interactúa con el sistema inteligente & La UI cambia las tareas, los datos, el negocio y la dirección de la investigación. \\
Different Bet & Una ruta de apuesta distinta a la del actor dominante & No es copiar con otra apariencia, sino cambiar las condiciones de frontera. \\
\bottomrule
\end{tabularx}
\end{knowledgebox}

\begin{knowledgebox}{Ruta de lectura del episodio}
\begin{tabularx}{\textwidth}{p{0.22\textwidth}p{0.32\textwidth}X}
\toprule
Ruta & Pregunta central & Secciones correspondientes \\
\midrule
Persona & ¿Por qué el lenguaje es una herramienta de generalización y por qué Yao Shunyu apuesta por Agent? & Lenguaje y posturas no consensuadas. \\
Sistema & ¿Cómo evolucionan juntas las líneas de tareas, entornos y métodos de Agent? & Las dos preguntas centrales y las tres oleadas de auge y declive. \\
Recompensa & ¿Cómo obtiene un Agent reward, autoexploración y estructura multi-agent? & Reward, code affordance, generalización. \\
Frontera & ¿Cómo determinan Chatbot, interface y Super App las nuevas oportunidades? & La absorción de fronteras y las formas de interacción. \\
Panorama global & ¿Cómo coexisten lo unipolar y lo plural, OpenAI y las oportunidades de emprendimiento? & El panorama humano global y different bet. \\
\bottomrule
\end{tabularx}
\end{knowledgebox}

\subsection{Resumen de la sección}

EP115 es una entrevista sobre la teoría de Agent, no solo una entrevista a un investigador de OpenAI. Reúne en un mismo mapa el lenguaje, las personas, las tareas, los entornos, las recompensas, la interacción y la competencia entre plataformas, por lo que sirve como nota central para entender la segunda mitad de los Agent.
